\documentclass[11pt,a4paper]{article}
\usepackage[T1]{fontenc}
\usepackage[utf8]{inputenc}
\usepackage[czech]{babel}
\usepackage{amsmath,amssymb,amsthm,mathtools}
\usepackage[margin=2.5cm]{geometry}
\usepackage{microtype}
\newtheorem{theorem}{Věta}
\newtheorem{lemma}[theorem]{Lemma}
\newtheorem{corollary}[theorem]{Důsledek}
\newcommand{\E}{\mathbb E}
\newcommand{\Prb}{\mathbb P}
\newcommand{\st}{\mathrm{st}}
\newcommand{\Law}{\operatorname{Law}}
\title{Orbitální uzávěr Waringových směsí:\\proč konečné splity nemohou pokračovat neomezeně}
\author{}
\date{19. září 2026}
\begin{document}
\maketitle

\noindent\textbf{Výsledek.}
Pro danou dvojici $P,Q$ neexistuje nekonečný strom dominancí
\[
Q*R_{\nu_h}\succeq_{\st}
\sum_{x=0}^4P_x\delta_x*R_{\nu_{hx}},\qquad
R_\nu=\int W_a\,\nu(da),
\]
jehož kořen má konečný fázový kapitál $\int a\,\nu_\varnothing(da)$.
To zahrnuje kořen $H_2=W_2$ a libovolné history-dependent,
spojité či nekonečně podporované fázové míry.
Předchozí explicitní dva stupně zůstávají platnými konečnými
certifikáty. Tato věta ukazuje, že čisté Waringovy směsi nelze
uzavřít do požadované nekonečné orbity.
Obecnou logaritmickou horní mez pro exactní stromy věta nevyvrací.

\section{Definice a dvě různé střední hodnoty}
Pracujeme s
\[
P=\tfrac1{64}(8,16,16,16,8),\qquad
Q=\tfrac1{64}(7,20,10,20,7).
\]
Platí $\E P=\E Q=2$ a $Q_0/P_0=7/8<1$.
Pro $a>0$ je
\[
\Prb(W_a\ge j)=\frac{a}{a+j-1},\qquad j\ge1.
\]
Pro potřeby uzavřenosti rodinu \emph{rozšíříme} o $W_0=\delta_1$.
Vyvrácení existence v této větší rodině implikuje vyvrácení
existence v původní rodině $a>0$.

Nechť $\nu$ je pravděpodobnostní míra na $[0,\infty)$. Označme
\[
R_\nu=\int W_a\,\nu(da),\qquad
m(\nu)=\int a\,\nu(da),\qquad
r(\nu)=R_\nu(1)=\int\frac{1}{1+a}\,\nu(da).
\]
$m(\nu)$ je střední hodnota \emph{fázového parametru}, nikoli
střední rezerva. Střední rezerva může být nekonečná.
Pro $k\ge1$ a $j\ge2$ máme
\begin{equation}\label{cdf}
F_{R_\nu}(k)=\int\frac{k}{a+k}\,\nu(da),\qquad
\Prb(R_\nu\ge j)=\int\frac{a}{a+j-1}\,\nu(da).
\end{equation}
Všechny $R_\nu$ jsou podporovány na $\{1,2,\ldots\}$.

Nekonečným referenčním stromem rozumíme pravděpodobnostní míry
$(\nu_h)_{h\in\{0,1,2,3,4\}^*}$ splňující
\begin{equation}\label{fan}
Q*R_{\nu_h}\succeq_{\st}
\sum_xP_x\delta_x*R_{\nu_{hx}}
\quad\text{pro každé konečné }h.
\end{equation}
Nejde o předpoklad rovnosti skutečných posteriorů s referencemi.
Vyšetřujeme právě navrženou postačující konstrukci \eqref{fan}.

\section{Kapitálový účet a kompaktnost}
\begin{lemma}[Automatický kapitálový účet]\label{capital}
Pokud rodič v \eqref{fan} má konečný fázový kapitál, pak
\begin{equation}\label{supermart}
\sum_xP_xm(\nu_{hx})\le m(\nu_h).
\end{equation}
Zvláště pro $p(h)=\prod_iP_{h_i}$ platí
$m(\nu_h)\le m(\nu_\varnothing)/p(h)$.
\end{lemma}
\begin{proof}
Z \eqref{cdf} plyne
\[
\lim_{j\to\infty}j\Prb(R_\nu\ge j)=m(\nu).
\]
Pro konečné $m$ použijeme dominovanou konvergenci s majorantou
$2a$ pro $j\ge2$. Pro nekonečné $m$ dává Fatouovo lemma nekonečný
liminf, tedy tutéž identitu v rozšířeném smyslu.
Konvoluce s mírou na $\{0,\ldots,4\}$ tento limit nemění.
Použití limitu na ocasovou nerovnost \eqref{fan} dokazuje
\eqref{supermart}, včetně konečnosti kapitálů všech dětí.
Indukce dává odhad pro historii $h$.
\end{proof}

\begin{lemma}[Uzavřenost při omezeném kořenovém kapitálu]\label{closed}
Nechť $\nu^{(i)}$ jsou kořeny nekonečných referenčních stromů
a $m(\nu^{(i)})\le C<\infty$. Existuje podposloupnost celých stromů
konvergující slabě v každém pevném uzlu k nekonečnému referenčnímu
stromu. Jeho kořenový kapitál je nejvýše
$\liminf_i m(\nu^{(i)})$, pokud podposloupnost vybíráme podél liminf.
\end{lemma}
\begin{proof}
Pro každé pevné $h$ platí podle lemmatu \ref{capital}
\[
\nu_h^{(i)}([A,\infty))\le\frac{C}{p(h)A}.
\]
Fázové míry v tomto uzlu jsou tedy těsné. Slabou kompaktností
pravděpodobnostních měr a diagonálním výběrem přes spočetně mnoho
historií dostaneme slabé limity $\nu_h$ současně ve všech uzlech.
První moment je zdola polospojitý, takže příslušné momentové meze
zůstávají platné.

Funkce $a\mapsto k/(a+k)$ je pro každé $k\ge1$ omezená a spojitá
na $[0,\infty)$. Každá konkrétní CDF nerovnost v \eqref{fan} proto
přejde do limity; pro $k\le0$ jsou CDF nulové.
Všechny limitní míry jsou pravděpodobnostní. Tím je ověřen celý
nekonečný strom.

Konstanty $C/p(h)$ závisejí na historii. V tomto kroku se
nepředpokládá jednotná těsnost přes všechny historie.
\end{proof}

\begin{lemma}[Minimální životaschopný kapitál]\label{minimum}
Pokud existuje alespoň jeden nekonečný referenční strom s konečným
kořenovým kapitálem, pak existuje takový strom a číslo $c<\infty$,
pro něž
\begin{equation}\label{constant}
m(\nu_h)=c\quad\text{pro všechny historie }h.
\end{equation}
\end{lemma}
\begin{proof}
Označme $\mathcal V$ množinu všech fázových měr, které jsou kořeny
nekonečných referenčních stromů s konečným kapitálem. Z předpokladu
je neprázdná. Položme
\[
c=\inf_{\nu\in\mathcal V}m(\nu)\in[0,\infty).
\]
Vybereme kořeny s kapitály klesajícími k $c$ a použijeme lemma
\ref{closed}. Limitní kořen stále patří do $\mathcal V$ a má kapitál
nejvýše $c$, tedy právě $c$.

Každé jeho dítě je kořenem vlastního nekonečného podstromu, a proto
má kapitál alespoň $c$. Nerovnost \eqref{supermart} a $P_x>0$ nutí
všechny tyto kapitály být právě $c$. Stejný argument v každém
potomku dokazuje \eqref{constant}.
\end{proof}

\section{Stacionární dominance nemůže ztrácet rezervu}
\begin{lemma}[Nulový stacionární slack]\label{stationary}
Nechť $J,Z$ jsou nezáporné celočíselné veličiny se stejným
marginálním zákonem, $0\le X,Y\le4$, $\E X=\E Y$, a
\[
\Law(J+Y)\succeq_{\st}\Law(Z+X).
\]
Pak ve skutečnosti $\Law(J+Y)=\Law(Z+X)$.
Nezávislost není potřebná a střední hodnoty $J,Z$ smějí být nekonečné.
\end{lemma}
\begin{proof}
Položme $\sigma(k)=F_{Z+X}(k)-F_{J+Y}(k)\ge0$.
Protože $\E\min(J,N)=\E\min(Z,N)$, máme
\begin{align*}
\sum_{k=0}^{N-1}\sigma(k)
&=\E\min(J+Y,N)-\E\min(Z+X,N)\\
&=\E[\min(J+Y,N)-\min(J,N)]\\
&\quad-\E[\min(Z+X,N)-\min(Z,N)].
\end{align*}
Oba rozdíly v hranatých závorkách leží v $[0,4]$ a konvergují
k $Y$, respektive $X$. Dominovanou konvergencí tedy pravá strana
konverguje k nule. Levá strana je rostoucí součet nezáporných
čísel, a proto každý člen $\sigma(k)$ musí být nulový.
\end{proof}

\section{Orbitální spor}
\begin{theorem}[Nemožnost nekonečné Waringovy orbity]\label{main}
Neexistuje nekonečný referenční strom \eqref{fan} s
$m(\nu_\varnothing)<\infty$. Zejména neexistuje takový strom
s kořenem $\nu_\varnothing=\delta_2$.
\end{theorem}
\begin{proof}
Předpokládejme existenci. Podle lemmatu \ref{minimum} můžeme
přejít k pomocnému nekonečnému stromu s konstantním kapitálem $c$.
Množina
\[
K_c=\left\{\nu\in\mathcal P([0,\infty)):
\int a\,\nu(da)\le c\right\}
\]
je slabě kompaktní: první moment dává těsnost a jeho dolní
polospojitost dává uzavřenost. Všechny fáze pomocného stromu
náležejí do $K_c$.

Pro $t\ge0$ zaveďme pravděpodobnostní míru na $K_c^6$
\[
\rho_t=\sum_{|h|=t}p(h)
\delta_{(\nu_h,\nu_{h0},\ldots,\nu_{h4})},\qquad
\pi_N=\frac1N\sum_{t=0}^{N-1}\rho_t.
\]
Vybereme slabě konvergentní podposloupnost $\pi_N\to\pi$.
Množina šestic splňujících \eqref{fan} je uzavřená, neboť každá
CDF nerovnost je spojitá v těchto šesti fázových mírách. Proto je
i $\pi$ nesena přípustnými fany.

Označme $\eta$ rodičovský marginál $\pi$ a $\eta_x$ marginál
dítěte $x$. Teleskopický součet přes hloubky dává
\begin{equation}\label{balance}
\boxed{\eta=\sum_xP_x\eta_x.}
\end{equation}
Přesněji, odpovídající rozdíl u $\pi_N$ je
\[
\frac1N\left(\delta_{\nu_\varnothing}
-\sum_{|h|=N}p(h)\delta_{\nu_h}\right),
\]
jehož totální variace je nejvýše $2/N$.

Položme $\overline R=\int R_\nu\,\eta(d\nu)$ a zprůměrujme
dominance přes $\pi$. Levou stranu realizujeme jako $J+Y$, kde
$J\sim\overline R$ a $Y\sim Q$ nezávisle. Pravou stranu
realizujeme výběrem šestic podle $\pi$, dále $X\sim P$ nezávisle
na šestici a $Z\sim R_{\nu_X}$. Z \eqref{balance} má $Z$ právě
zákon $\overline R$, tedy stejný jako $J$.
Protože $\E X=\E Y=2$, lemma \ref{stationary} ukazuje, že
zprůměrovaná dominance je rovnost.

Lokální CDF slack je nezáporný v každé šestici. Jeho integrál je
v každém prahu nulový, takže slack je nulový $\pi$-skoro jistě.
V prahu $k=1$ to znamená
\begin{equation}\label{head}
P_0r(\nu_0)=Q_0r(\nu),\qquad
r(\nu_0)=\frac78r(\nu)
\quad\pi\text{-skoro jistě}.
\end{equation}
Zde jsme použili, že všechny reference jsou podporovány na
kladných celých číslech.

Jensenova nerovnost na $K_c$ dává
\[
r(\nu)=\int\frac1{1+a}\,\nu(da)
\ge\frac1{1+m(\nu)}\ge\frac1{1+c}>0.
\]
Nechť $\ell=\operatorname*{ess\,inf}_{\eta}r$; tedy
$\ell\ge1/(1+c)>0$.
Z \eqref{balance} a $P_0>0$ plyne $\eta_0\ll\eta$, a tudíž
$r(\nu_0)\ge\ell$ pro $\pi$-skoro každou šestici.
Rovnice \eqref{head} nyní nutí $r(\nu)\ge8\ell/7$ rodičovsky
skoro jistě. Z definice esenciálního infima tedy
\[
\ell\ge\frac87\ell,
\]
což odporuje $\ell>0$.
\end{proof}

\begin{corollary}[Konečné horizonty nemohou přetrvat neomezeně]
Existuje konečné $N$, pro které z $\nu_\varnothing=\delta_2$
neexistuje referenční strom \eqref{fan} hloubky $N$.
Tento důkaz nedává explicitní numerickou hodnotu $N$.
\end{corollary}
\begin{proof}
Kdyby existovaly stromy pro neomezené horizonty, v každém pevném
uzlu by jejich fázové první momenty byly nejvýše $2/p(h)$.
Stejný diagonální argument jako v lemmatu \ref{closed} by dal
nekonečný strom s kořenem $\delta_2$, ve sporu s větou \ref{main}.
Kompatibilita původních konečných svědků není potřebná.
\end{proof}

\section{Co bylo odvozeno a co se nepředpokládalo}
\begin{itemize}
\item Žádná horní mez parametrů $a$, počtu fází ani podpory
fázových měr není předpokládána.
\item Ve výchozí hypotéze nejsou kapitály jednotlivých historií
jednotně omezené. Automatický odhad $m_h\le m_\varnothing/p(h)$
smí s hloubkou růst.
\item Jednotná kompaktnost vzniká až v pomocném stromu dosažením
minimálního životaschopného kapitálu. Je dokázaným důsledkem
hypotetické existence, nikoli dodatečným požadavkem na konstrukci.
\item Střední rezervy nejsou nikde prohlášeny za konečné.
Stacionární argument používá výhradně ořezané rezervy a přírůstky
omezené číslem $4$.
\item Výsledek se týká dominancí mezi Waringovými referencemi.
Neříká, že libovolný skutečný exactní strom musí mít takové reference.
\end{itemize}

\section{Stav potřebný pro další konstrukci}
U čistých směsí platí pro $j\ge2$
\[
\Prb(R_\nu\ge j)\le\frac{m(\nu)}{j-1}.
\]
Tentýž skalární kapitál tedy omezuje ocasový koeficient i všechny
kvantily. Tato vlastnost umožnila minimální kompaktní orbitu a spor.
Samotné přesouvání parametrů $a$ do nekonečna ji neodstraní.

Přirozené širší reference mají například tvar
\[
R_\lambda=\int\delta_r*W_a\,\lambda(dr,da),\qquad r\in\mathbb N_0.
\]
Již rodina $\delta_r*W_2$ má stále ocasový koeficient $2$, ale při
$r\to\infty$ není těsná. Proto se na tuto širší rodinu uvedený
argument minimálního kapitálu nepřenáší. Pro takový posunutý stav
však zde není sestrojen orbitální splitting operator.

\medskip
\noindent\textbf{Závěr.}
Čistá Waringova směs není dostatečný stav pro zamýšlené orbitální
uzavření. Explicitní dva stupně byly korektní, ale tento konkrétní
mechanismus logaritmus uzavřít nemůže. Obecná konstrukce
$\mathcal B_n=O(\log n)$ zůstává tímto výsledkem nevyřešena.
\end{document}
